
\documentclass[11pt,a4paper,sans]{moderncv}        
\moderncvstyle{classic}                             
\moderncvcolor{blue}                               
\nopagenumbers{}                                
\usepackage[scale=0.8]{geometry}
\usepackage{amssymb, latexsym, amsfonts, pigpen, enumitem, mathrsfs, amsthm, bbm, etaremune}
\setlength{\hintscolumnwidth}{3.65cm}                % if you want to change the width of the column with the dates
%\setlength{\makecvtitlenamewidth}{10cm}           % for the 'classic' style, if you want to force the width allocated to your name and avoid line breaks. be careful though, the length is normally calculated to avoid any overlap with your personal info; use this at your own typographical risks...

\usepackage{xcolor}
\usepackage{hyperref}
\hypersetup{colorlinks=true,linkcolor=blue,urlcolor=blue}

% personal data
\name{Alexey}{Ananyevskiy}
\title{Curriculum Vitae, September 2026}                               % optional, remove / comment the line if not wanted
%\address{Chebyshev Laboratory, SPbU, 14th Line 29B}{199178, St. Petersburg, Russia}% optional, remove / comment the line if not wanted; the "postcode city" and "country" arguments can be omitted or provided empty
%\phone[mobile]{+79119321338}                   % optional, remove / comment the line if not wanted; the optional "type" of the phone can be "mobile" (default), "fixed" or "fax"
\email{alseang@gmail.com}                               % optional, remove / comment the line if not wanted
\homepage{aananyevskiy.github.io}                          % optional, remove / comment the line if not wanted

% bibliography adjustements (only useful if you make citations in your resume, or print a list of publications using BibTeX)
%   to show numerical labels in the bibliography (default is to show no labels)
\makeatletter\renewcommand*{\bibliographyitemlabel}{\@biblabel{\arabic{enumiv}}}\makeatother
%   to redefine the bibliography heading string ("Publications")
%\renewcommand{\refname}{Articles}

% bibliography with mutiple entries
%\usepackage{multibib}
%\newcites{book,misc}{{Books},{Others}}
%----------------------------------------------------------------------------------
%            content
%----------------------------------------------------------------------------------
\begin{document}
%\begin{CJK*}{UTF8}{gbsn}                          % to typeset your resume in Chinese using CJK
%-----       resume       ---------------------------------------------------------
\makecvtitle


\section{Employment}
\cventry{2022--\hphantom{2017}}{Heisenberg Position}{}{}{Ludwig-Maximilians-Universität München, Germany}{}
\cventry{2017--2022}{Deputy Director}{}{}{EIMI, St. Petersburg, Russia}{}
\cventry{2017--2022}{Senior Researcher}{}{}{PDMI RAS, St. Petersburg, Russia}{}
\cventry{2017--2019}{Senior Researcher}{}{}{Chebyshev Laboratory at SPbU, St. Petersburg, Russia}{}
\cventry{2015--2019}{Deputy Head}{}{}{Chebyshev Laboratory at SPbU, St. Petersburg, Russia}{}
\cventry{2013--2017}{Postdoc}{}{}{Chebyshev Laboratory at SPbU, St. Petersburg, Russia}{}
\cventry{2010--2013}{Lecturer}{}{}{Algebra Department at SPbU, St. Petersburg, Russia}{}

\section{Education}
\cventry{2010--2013}{Ph.D. in Mathematics}{SPbU, St. Petersburg, Russia}{}{}{}
\cventry{2005--2010}{Specialist in Mathematics}{SPbU, St. Petersburg, Russia}{}{}{}

%\section{Employment}
%\cvitem{title}{\emph{Title}}
%\cvitem{supervisors}{Supervisors}
%\cvitem{description}{Short thesis abstract}

\section{Research visits}
\cvitem{Oct -- Nov 2026}{SLMath, Berkeley, USA (planned)}
\cvitem{Mar 2020}{Isaac Newton Institute, Cambridge, UK; suspended due to the COVID-19 pandemic}
\cvitem{Oct 2018}{University of Oslo, Oslo, Norway}
\cvitem{Feb -- Mar 2018}{University of Oslo, Oslo, Norway}
\cvitem{Jan -- Apr 2017}{Mittag-Leffler Institute, Stockholm, Sweden}
\cvitem{Nov 2016}{Universit\"at Osnabr\"uck, Osnabr\"uck, Germany}
\cvitem{Jan -- Apr 2015}{Institute for Advanced Study, Princeton, USA}
\cvitem{Jun 2014}{Universit\"at Duisburg-Essen, Essen, Germany}
\cvitem{May 2013}{Fields Institute, Toronto, Canada}
\cvitem{Jul -- Aug 2011}{Universit\"at Bielefeld, Bielefeld, Germany}
\cvitem{Jun -- Jul 2009}{Universit\"at Bielefeld, Bielefeld, Germany}

\section{Research interests}
motivic homotopy theory, vector bundles, algebraic $\mathrm{K}$-theory, homogeneous varieties, algebraic geometry

\section{Grants, Fellowships and Awards}
\subsection{Grants}
\cvitem{2023--\hphantom{2022}}{DFG research grant AN 1545/4-1, project "Equivariant and weak orientations in the motivic homotopy theory", 226.650 euro. Joint with Nikita Geldhauser.}
\cvitem{2022--\hphantom{2022}}{DFG Heisenberg Program, grant AN 1545/1-1, project "Motivic homotopy theory and cohomological properties of algebraic varieties", 540.000 euro.}
\cvitem{2021--2022}{"Basis" Foundation grant "Stable motivic homotopy theory, K-theory and generalized motivic cohomology", Principal Investigator, 1.200.000 RUB/year.}
\cvitem{2020--2022}{joint RSF-DFG grant 20-41-04401 "Algebraic bordism spectra: Computations, filtrations, applications", Principal Investigator on the Russian side, 6.000.000 RUB/year on the Russian side. Principal Investigator on the German side is Oliver R\"ondigs.}
\cvitem{2014--2015}{RFBR grant 14-01-31095-mol\_a \emph{"Homotopical methods in algebraic geometry and derived Witt groups"}, Principal Investigator, 400.000 RUB/year.}
\cvitem{Investigator\\ (not PI)}{RFBR grants 19-01-00513-a, 18-31-20044-mol\_a\_ved, 16-01-00750-a, 15-01-03034-a, 13-01-00429-a, 12-01-33057-mol\_a\_ved, 12-01-31100-mol\_a, 10-01-00551\_a, RSF grant 14-11-00456}
\subsection{Fellowships and Awards}
\cvitem{2020}{Simons Foundation Fellowship for a stay at the Isaac Newton Institute, Cambridge}
\cvitem{2018--2020}{"Young Russian Mathematics" fellowship}
\cvitem{2017}{Fellowship for a stay at Mittag-Leffler Institute, Stockholm}
\cvitem{2015}{IAS fellowship to stay at Institute for Advanced Study, Princeton}
\cvitem{2013--2015}{Dynasty foundation fellowship for Russia's young mathematicians}
\cvitem{2013, 2015}{Gazprom Neft prize for young mathematicians}
\cvitem{2012}{Young Mathematician prize of the St. Petersburg Mathematical Society}
\cvitem{2010, 2012}{Rokhlin prize for St. Petersburg young mathematicians}

\section{Publications}
\begin{etaremune}[leftmargin=20pt]
	\item
	Alexey Ananyevskiy, Alexander Samokhin, \emph{Semiorthogonal decompositions for families of twisted flag varieties}, \url{https://arxiv.org/abs/2609.05203}
	\item
	Alexey Ananyevskiy, Nikita Geldhauser, \emph{Chow rings of quasi-split geometrically almost simple algebraic groups}, to appear in Journal für die reine und angewandte Mathematik, \url{https://arxiv.org/abs/2408.09390}
	\item
	Alexey Ananyevskiy, Marc Levine, \emph{Combing a hedgehog over a field}, Algebra and Number Theory 19:12 (2025), 2409--2431, \url{https://doi.org/10.2140/ant.2025.19.2409}	 
	\item
	Alexey Ananyevskiy, Elden Elmanto, Oliver R\"ondigs, Maria Yakerson, \emph{On the motivic Adams conjecture}, Journal of Topology 18:2 (2025), e70026, \url{https://doi.org/10.1112/topo.70026}
	\item
	Alexey Ananyevskiy, \emph{On the $\mathbb{A}^1$-Euler characteristic of the variety of maximal tori in a reductive group}, IMRN, 2022:19 (2022), 15390–15409, \url{https://doi.org/10.1093/imrn/rnab156}
	\item
	Alexey Ananyevskiy, \emph{Thom isomorphisms in triangulated motivic categories}, Algebraic and Geometric Topology, 21:4 (2021), 2085-2106, \url{https://doi.org/10.2140/agt.2021.21.2085}
	\item
	Alexey Ananyevskiy, Grigory Garkusha, Ivan Panin, \emph{Cancellation theorem for framed motives of algebraic varieties}, Advances in Mathematics, 383 (2021), 107681, \url{http://doi.org/10.1016/j.aim.2021.107681}
	\item
	Alexey Ananyevskiy, \emph{SL-oriented cohomology theories}, in \emph{Motivic Homotopy Theory and Refined Enumerative Geometry}, Contemporary Mathematics, 745 (2020), 1-20, \url{https://doi.org/10.1090/conm/745}
	\item	
	Alexey Ananyevskiy, Oliver R\"ondigs, Paul Arne \O stv\ae r, \emph{On very effective hermitian $K$-theory}, Mathematische Zeitschrift, 294 (2020), 1021-1034, \url{https://doi.org/10.1007/s00209-019-02302-z}
	\item
	Alexey Ananyevskiy, Alexander Neshitov, \emph{Framed and MW-transfers for homotopy modules}, Selecta Mathematica, 25:26 (2019), \url{https://doi.org/10.1007/s00029-019-0472-0}
	\item	
	Alexey Ananyevskiy, Andrei Druzhinin, \emph{Rigidity for linear framed presheaves and generalized motivic cohomology theories}, Advances in Mathematics, 333 (2018), 423-462, \url{https://doi.org/10.1016/j.aim.2018.05.013}
	\item
	Alexey Ananyevskiy, \emph{On the zeroth stable $\mathbb{A}^1$-homotopy group of a smooth curve}, Journal of Pure and Applied Algebra, 222:10 (2018), 3195-3218, \url{https://doi.org/10.1016/j.jpaa.2017.12.001}
	\item	
	Alexey Ananyevskiy, Marc Levine, Ivan Panin, \emph{Witt sheaves and the $\eta$-inverted sphere spectrum}, Journal of Topology, 10:2 (2017), 370-385, \url{https://doi.org/10.1112/topo.12015}
	\item	
	Alexey Ananyevskiy, \emph{Stable operations and cooperations in derived Witt theory with rational coefficients}, Annals of K-theory, 2:4 (2017), 517-560, \url{https://doi.org/10.2140/akt.2017.2.517}
	\item
	Alexey Ananyevskiy, \emph{On the zeroth stable $\mathbb{A}^1$-homotopy group of a smooth projective variety}, Journal of Mathematical Sciences, 222:4 (2017), 367-369 (translated from Russian), \url{https://doi.org/10.1007/s10958-017-3306-7}
	\item	
	Alexey Ananyevskiy, \emph{On the push-forwards for motivic cohomology theories with invertible stable Hopf element}, Manuscripta Mathematica, 150:1-2 (2016), 21-44, \url{https://doi.org/10.1007/s00229-015-0799-6}
	\item
	Alexey Ananyevskiy, \emph{On the relation of special linear algebraic cobordism to Witt groups}, Homology, Homotopy and Applications, 18:1 (2016), 205-230, \url{https://doi.org/10.4310/HHA.2016.v18.n1.a11}
	\item
	Alexey Ananyevskiy, \emph{The special linear version of the projective bundle theorem}, Compositio Mathematica, 151:3 (2015), 461-501, \url{https://doi.org/10.1112/S0010437X14007702}
	\item
	Alexey Ananyevskiy, Asher Auel, Skip Garibaldi, Kirill~Zainoulline, \emph{Exceptional collections of line bundles on projective homogeneous varieties}, Advances in Mathematics, 236 (2013), 111-130, \url{https://doi.org/10.1016/j.aim.2012.12.016}
	\item
	Alexey Ananyevskiy, \emph{Relationship between algebraic MSL-cobordisms and derived Witt groups}, Doklady Mathematics, 87:1 (2013), 76-78, \url{http://doi.org/10.1134/S1064562413010286}
	\item
	Alexey Ananyevskiy, \emph{On the algebraic $K$-theory of some homogeneous varieties}, Documenta Mathematica, 17 (2012), 167-193, \url{https://www.emis.de/journals/DMJDMV/vol-17/07.html}
	\item
	Alexey Ananyevskiy, Nikolai Vavilov, Sergey Sinchuk, \emph{Overgroups of $E(m;R)\otimes E(n;R)$. I. Levels and normalisers,} St. Petersburg Mathematical Journal, 23:5 (2012), 819-849 (translated from Russian), \url{http://doi.org/10.1090/S1061-0022-2012-01219-7}
	\item
	Alexey Ananyevskiy, Nikolai Vavilov, Sergey Sinchuk, \emph{On the description of overgroups of $E(m;R)\otimes E(n;R)$,} Journal of Mathematical Sciences, 161:4 (2009), 461-473 (translated from Russian), \url{https://doi.org/10.1007/s10958-009-9576-y}
\end{etaremune}

\section{Selected talks (last five years)}
\cvitem{September 10, 2026}{\emph{Combing hedgehogs over a field}, conference \emph{Geometria in Bicocca 2026}, University of Milano-Bicocca, Milan, Italy}
\cvitem{September 08, 2026}{\emph{Combing hedgehogs over a field}, DMV Annual Meeting, Konstanz University, Konstanz, Germany}
\cvitem{December 05, 2025}{\emph{On semiorthogonal decompositions for twisted flag varieties},	\emph{Complex algebraic geometry seminar}, Università degli Studi di Milano, Milan, Italy}
\cvitem{September 04, 2025}{\emph{On the semiorthogonal decompositions for twisted flag varieties}, Workshop \emph{Semiorthogonal decompositions for representations of algebraic groups}, Universität Bielefeld, Bielefeld, Germany}
\cvitem{March 17, 2025}{\emph{Nowhere vanishing sections of vector bundles revisited}, Conference \emph{Motivic homotopy theory}, Universität Regensburg, Regensburg, Germany}
\cvitem{October 22, 2024}{\emph{Combing a hedgehog over a field}, \emph{Algebraic geometry seminar}, University of Zurich, Zurich, Switzerland}
\cvitem{January 15, 2024}{\emph{Combing a hedgehog over a field}, \emph{Algebra seminar}, Charles University, Prague, Czech Republic}
\cvitem{November 23, 2023}{\emph{Combing a hedgehog over a field}, \emph{Kolloquium Geometrie und Arithmetik}, Johannes Gutenberg Universität Mainz, Mainz, Germany}
\cvitem{October 19, 2023}{\emph{On the motivic Adams conjecture}, \emph{AG seminar}, Universität Regensburg, Regensburg, Germany}
\cvitem{September 28, 2023}{\emph{Non-Vanishing Sections of Vector Bundles and Chern Classes}, \emph{ESAGA seminar}, Universität Duisburg--Essen, Essen, Germany}
\cvitem{June 22, 2023}{\emph{On the existence of non-vanishing sections of vector bundles}, \emph{Séminaire d'arithmétique à Lyon}, ENS de Lyon, Lyon, France}
\cvitem{May 10, 2023}{\emph{Non-Vanishing Sections of Algebraic Vector Bundles and Trivial Chern Classes}, \emph{Colloquium of Institut für Mathematik}, Universität Osnabrück, Germany}
\cvitem{November 17, 2022}{\emph{On the existence of a nowhere vanishing section for a vector bundle with trivial top Chern class}, \emph{Oberseminar zur Algebra und Zahlentheorie}, Universität Augsburg, Augsburg, Germany}
\cvitem{October 27, 2022}{\emph{On the existence of a nowhere vanishing section for a vector bundle with trivial top Chern class}, seminar \emph{Motivic algebraic topology}, LMU, Munich, Germany}




\section{Mentorship}
\cvitem{2026}{Andriy Kryzhanovskyy, MSc student at LMU Munich}
\cvitem{2025--2026}{Anton Lin, MSc student at LMU Munich}
\cvitem{2024--\hphantom{2020}}{Egor Zolotarev, PhD student at LMU Munich}
\cvitem{2021--2022}{Tariq Syed, Postdoc at EIMI}

%\section{Teaching philosophy}
%\parindent=15pt
%One of the most important things in teaching is to adapt to the audience. Students with distinct background and at different stages of education may need completely different approaches, and teaching has to pursue different goals besides simply conveying some mathematical knowledge. From my point of view, the main goal for teaching the first year students, especially those who have had only a very limited experience with advanced mathematics, is to explain what does it mean to prove a statement and to make students comfortable with mathematical rigor and basic abstract constructions. At this stage it is very important to provide all the details and give clear expositions, since these first proofs that the students encounter will serve as role model for their future reasoning. Then, at the second and third years it is very important to show the rich diversity of mathematical research and provide connections between different disciplines within the field so that one does not think that mathematics is just a collection of tricks. By the final year of Bachelor's studies the students usually become more self-sustained and some technical details in the expositions could be omitted and left as exercises, but one should always clearly state the claims whose proofs are omitted and all the statements should be completely unambiguous. Masters students usually attend more specialized courses, but it remains to be important to keep connections between different branches of mathematics. Also at this stage some material should be learned at a reading seminar or given as a supplementary reading for the lecture course, introducing the students to the practices of the actual research.
%\\
%\indent
%At all levels abstract constructions should be supplied with rich collections of particular examples, algebraic manipulations should be accompanied with geometric pictures whenever it is possible, and theorems and propositions should be followed by examples and counter-examples demonstrating the necessity of assumptions. Another crucial thing is to maintain connection with the audience and to adapt the exposition based on the feedback, slowing down and providing more details if necessary or leaving more technicalities as exercises if the audience is really engaged.
%\\
%\indent
%The above addressed mostly teaching classes, while supervision of individual students and direction of theses is another thing. In this case the choice of the topic and the problem is the key, it should not be too easy for the student but it should be possible to deal with. Again, here a constant contact is of uttermost importance, one of the worst things that the supervisor can do is to give a problem and say to come back with the advances, this usually leads that the student disappears for a long time and eventually comes back depressed and without any success. One has to arrange regular meetings for discussions, this will keep the student motivated and going. It is also quite important to develop a relationship of trust with the supervised student, the situation where the student hesitates to admit that they are struggling with the problem is highly undesirable. In general, I think that supervised students need a lot of support and guidance.

\section{Services to the community}
\cvitem{Refereeing}{Adv. Math.,  Algebr. Geom. Topol., Ann. Inst. Fourier, Ann. K-theory, Compos. Math., Doc. Math., IMRN, Invent. Math., Izv. Math., J. Algebra, J. Inst. Math. Jussieu, J. Pure Appl. Algebra, J. Topology, Math. Proc. Cambridge Philos. Soc., Math Z., Pacific J. Math., St. Petersburg Math. J., Transform. Groups, MathSciNet}
\cvitem{Reviewer}{German Research Foundation}
\cvitem{Organizing Zoom seminar}{Zoom seminar on $\mathbb{A}^1$-topology, motives and $\mathrm{K}$-theory in 2020-2022\newline \url{https://indico.eimi.ru/category/12/}.}
\cvitem{Organizing conferences}{
	\href{https://www.cas.lmu.de/en/events/event/motifs-and-motives-d8fb57b4.html}{\textbf{Motifs and Motives}}, June 29 -- July 03, 2026 at LMU, Munich, Germany. Organizers: Alexey Ananyevskiy, Nikita Geldhauser, Andrei Lavrenov and Maksim Zhykhovich.
	\newline
	\href{https://sites.google.com/view/agmspb2019}{\textbf{Emerging Research in Algebraic Groups, Motives, and K-Theory}}, September 9-13, 2019 at EIMI, St. Petersburg, Russia. Organizers: Alexey Ananyevskiy, Anastasia Stavrova and Nikolai Vavilov.
	\newline
	\href{https://sites.google.com/view/agmspb2019}{\textbf{Petersburg motives (dedicated to Ivan	Panin's 60th birthday)}}, September 2-6, 2019 at EIMI, St. Petersburg, Russia. Organizers: Alexey Ananyevskiy, Alexander Merkurjev and Serge Yagunov.
	\newline
	\href{https://sites.google.com/view/mispb2018}{\textbf{Motives in St. Petersburg 2018}}, September 3-14, 2018 at EIMI, St. Petersburg, Russia. Organizers: Alexey Ananyevskiy, Grigory Garkusha and Ivan Panin.
	\newline
	\href{https://chebyshev.spbu.ru/en/ag2018/}{\textbf{Moscow -- St. Petersburg algebraic geometry symposium for young mathematicians}}, May 10-11, 2018 at Chebyshev Laboratory, SPbU, St. Petersburg, Russia. Organizers: Alexey Ananyevskiy, Valery Gritsenko, Peter Zograf.
}


\end{document}

